\SourceLeaf{417}%
gens $aa'$ und in Beziehung auf denselben seien $\lambda,\Lambda$ dasselbe, was $\alpha,A$ und $\alpha',A'$ in Beziehung auf $a$ und $a'$ sind. Ist nun $\alpha\lambda$ oder $\Lambda-A$ dem Stück des Bogens $al$ proportional, so werden wir sagen, dass $aa'$ in seiner ganzen Länge gleichförmig gekrümmt sei und
\[
\frac{\Lambda-A}{al}
\]
das Krümmungsmass oder schlechtweg die Krümmung nennen; man sieht leicht, dass diess nur dann eintrifft, wenn $aa'$ wirklich ein Kreisbogen ist, und dass dann nach unserer Definition seine Krümmung $=\pm\dfrac1r$ sein wird, wenn $r$ den Halbmesser bedeutet. In so fern man immer $r$ als positiv ansieht, wird das obere oder untere Zeichen gelten, je nachdem das Centrum rechts oder links vom Bogen $aa'$ liegt ($a$ als Anfangs-, $a'$ als Endpunkt betrachtet, und die Richtungen im Hülfskreise von der Linken zur Rechten gezählt); die Veränderung Einer dieser Bedingungen ändert das Zeichen, die Veränderung von zweien stellt es wieder her.

Ist hingegen $\Lambda-A$ dem $al$ nicht proportional, so nennen wir den Bogen ungleichförmig gekrümmt und der Quotient
\[
\frac{\Lambda-A}{al}
\]
kann dann seine mittlere Krümmung heissen. Krümmung geradezu hingegen setzt immer die Bestimmung der Stelle voraus und wird durch die mittlere Krümmung eines Elements an dieser Stelle erklärt, ist also
\[
=\frac{d\Lambda}{dal}.
\]
Man sieht also, dass Bogen, Amplitudo und Krümmung in einem solchen Verhältniss zu einander stehen, wie Zeit, Bewegung und Geschwindigkeit, oder Volumen, Masse und Dichtigkeit. Das Reciprokum der Krümmung, nemlich
\[
\frac{dal}{d\Lambda}
\]
heisst der Krümmungshalbmesser an der Stelle $l$, und nach obigen Voraussetzungen heisst die Curve an dieser Stelle concav nach der Rechten, convex nach der Linken, wenn der Werth der Krümmung oder des Krümmungshalbmessers positiv ausfällt; bei einem negativen verhält es sich umgekehrt.
